\paragrafos{Πολλαπλασιασμός ρητών: το μείον επί μείον}{Α.7.5, «Πολλαπλασιασμός ρητών αριθμών»}

\begin{anagke}
Το βράδυ στο Μέτσοβο η θερμοκρασία πέφτει $2$ βαθμούς κάθε ώρα. Η Πολυξένη παρατηρεί το θερμόμετρο στις $20{:}00$, που δείχνει $0$°. «Τι θα δείχνει σε τρεις ώρες;». Ο Νίκος σκέφτεται: «Κάθε ώρα $-2$, τρεις ώρες: $-2-2-2=-6$». Η Αγγελική το γράφει ως πολλαπλασιασμό: $3\cdot(-2)=-6$. «Και πριν από τρεις ώρες, στις $17{:}00$;» ρωτάει η Μυρτώ. «Αφού τώρα πέφτει δύο βαθμούς ανά ώρα, πριν από τρεις ώρες είχε $6$ βαθμούς περισσότερο», απαντάει η Πολυξένη. Ο ρυθμός είναι $-2$ και ο χρόνος «τρεις ώρες πριν» είναι $-3$, και το αποτέλεσμα είναι $+6$. Άρα $(-2)\cdot(-3)=+6$.

Ο Στέλιος ξαφνιάζεται: «Το μείον επί μείον κάνει συν; Σαν τη διπλή άρνηση της γλώσσας που καταλήγει σε κατάφαση;». Η κυρία Αλεξάνδρα του απαντά: «Η γλώσσα είναι ωραία αναλογία, αλλά τα μαθηματικά θέλουν να δουν γιατί». Το ίδιο το μοτίβο δείχνει το γιατί: $3\cdot(-2)=-6$, $2\cdot(-2)=-4$, $1\cdot(-2)=-2$, $0\cdot(-2)=0$. Κάθε φορά που το πρώτο μικραίνει κατά $1$ το γινόμενο μεγαλώνει κατά $2$. Συνεχίζοντας: $(-1)\cdot(-2)=2$, $(-2)\cdot(-2)=4$.
\end{anagke}

\begin{lexi}{παράγοντας, γινόμενο, ομόσημοι}{παράγοντας, γινόμενο}
\lexiword{Παράγοντες} είναι οι αριθμοί που πολλαπλασιάζονται και \lexiword{γινόμενο} αυτό που «γίνεται». Η \lexiword{επιμεριστική} ιδιότητα είναι αυτή που «επιμερίζει» (μοιράζει σε μέρη) τον πολλαπλασιασμό πάνω σε ένα άθροισμα.
\end{lexi}

\begin{metafrasi}
Πέφτει $2$ βαθμούς την ώρα, για $3$ ώρες. & $3\cdot(-2)=-6$\\
Πριν από $3$ ώρες. & $(-2)\cdot(-3)=+6$\\
Ετερόσημοι. & γινόμενο αρνητικό\\
Ομόσημοι. & γινόμενο θετικό\\
Ο πολλαπλασιασμός πάνω σε άθροισμα. & $\alpha(\beta+\gamma)=\alpha\beta+\alpha\gamma$\\
\end{metafrasi}

\tmima{Η θεωρία}

\begin{kanonas}{γινόμενο ρητών}
Το γινόμενο δύο ρητών έχει απόλυτη τιμή το γινόμενο των απόλυτων τιμών τους και:
\begin{itemize}[leftmargin=1.2em,itemsep=0.3ex,topsep=2pt]
\item πρόσημο \textbf{θετικό} αν οι παράγοντες είναι ομόσημοι: $(-7)\cdot(-8)=+56$, $(+5)(+4)=+20$.
\item πρόσημο \textbf{αρνητικό} αν οι παράγοντες είναι ετερόσημοι: $(-6)\cdot(+4)=-24$.
\item Αν ένας παράγοντας είναι $0$, το γινόμενο είναι $0$.
\end{itemize}
\end{kanonas}

Για παράδειγμα $\left(-\frac{3}{4}\right)\cdot\left(-\frac{8}{9}\right)=+\frac{3\cdot8}{4\cdot9}=\frac{2}{3}$ και $(+0{,}5)\cdot(-12)=-6$.

\begin{idiotita}{ιδιότητες του πολλαπλασιασμού}
Για κάθε τρεις ρητούς ισχύουν: \textbf{αντιμεταθετική} $\alpha\beta=\beta\alpha$, \textbf{προσεταιριστική} $(\alpha\beta)\gamma=\alpha(\beta\gamma)$, $\alpha\cdot1=\alpha$, $\alpha\cdot0=0$, $\alpha\cdot(-1)=-\alpha$ και \textbf{επιμεριστική} $\alpha(\beta+\gamma)=\alpha\beta+\alpha\gamma$ (και $\alpha(\beta-\gamma)=\alpha\beta-\alpha\gamma$).
\end{idiotita}

\tmimasmall{Γινόμενο πολλών παραγόντων}

Το γινόμενο πολλών ρητών (όλων μη μηδενικών) είναι θετικό αν το πλήθος των αρνητικών παραγόντων είναι \emph{άρτιο} και αρνητικό αν είναι \emph{περιττό}. Η απόλυτη τιμή του είναι το γινόμενο των απόλυτων τιμών. Έτσι $(-2)\cdot3\cdot(-5)\cdot(-1)=-30$ (τρεις αρνητικοί).

Η επιμεριστική ιδιότητα δίνει συντομίες: $99\cdot(-6)=(100-1)\cdot(-6)=-600+6=-594$.

\begin{prosoxi}
\begin{itemize}[leftmargin=1.2em,itemsep=0.3ex,topsep=2pt]
\item Στον πολλαπλασιασμό, δύο αρνητικοί δίνουν θετικό, ενώ στην πρόσθεση δύο αρνητικοί δίνουν αρνητικό: $(-3)\cdot(-4)=+12$ και $(-3)+(-4)=-7$.
\item Το πρόσημο του γινομένου μετράει το πλήθος των αρνητικών παραγόντων, όχι το πλήθος των θετικών.
\item Στην επιμεριστική πολλαπλασιάζω τον αριθμό με \emph{κάθε} όρο της παρένθεσης: $(-2)(5-9)=-10+18$.
\item $(-3)^2$ και $-3^2$ δεν είναι ίδια, αλλά θα το δούμε στις δυνάμεις.
\end{itemize}
\end{prosoxi}

\tmima{Λυμένα παραδείγματα}

\paradeigma{Το μοτίβο}
\tixero Τα γινόμενα $3\cdot(-4)$, $2\cdot(-4)$, $1\cdot(-4)$, $0\cdot(-4)$.\par
\tithelo Να συνεχίσω το μοτίβο.\par
\possindeo Τα γινόμενα είναι $-12$, $-8$, $-4$, $0$ και αυξάνουν κατά $4$ κάθε φορά που ο πρώτος παράγοντας μειώνεται κατά $1$. Άρα $(-1)\cdot(-4)=4$, $(-2)\cdot(-4)=8$, $(-3)\cdot(-4)=12$.\par
\apantisi $4$, $8$, $12$.\par
\elegxos Ο κανόνας ομόσημοι δίνει θετικό: $(-3)\cdot(-4)=+12$.

\paradeigma{Γινόμενα με τον κανόνα}
\tixero Τα $(-6)(+4)$, $(-7)(-8)$, $(+0{,}5)(-12)$, $\left(-\frac{3}{4}\right)\left(-\frac{8}{9}\right)$.\par
\tithelo Να τα υπολογίσω.\par
\possindeo Πρώτα το πρόσημο, μετά οι απόλυτες τιμές. Ετερόσημοι: $-24$. Ομόσημοι: $+56$. Ετερόσημοι: $-6$. Ομόσημοι: $+\frac{24}{36}=\frac{2}{3}$.\par
\apantisi $-24$, $56$, $-6$, $\frac{2}{3}$.\par
\elegxos Το $7\cdot8=56$ και $6\cdot4=24$.

\paradeigma{Πολλοί παράγοντες}
\tixero Το γινόμενο $(-2)\cdot3\cdot(-5)\cdot(-1)$.\par
\tithelo Να βρω το πρόσημο και την τιμή.\par
\possindeo Αρνητικοί παράγοντες: $3$ (περιττό πλήθος), άρα αρνητικό. Απόλυτη τιμή: $2\cdot3\cdot5\cdot1=30$.\par
\apantisi $-30$.\par
\elegxos Βήμα βήμα: $(-2)\cdot3=-6$, $-6\cdot(-5)=30$, $30\cdot(-1)=-30$.

\paradeigma{Η επιμεριστική}
\tixero Οι $(-4)\cdot(7+3)$ και $99\cdot(-6)$.\par
\tithelo Να τις υπολογίσω εύκολα.\par
\possindeo $(-4)\cdot7+(-4)\cdot3=-28-12=-40$ (ή $-4\cdot10$). Για το δεύτερο: $99=100-1$, $(100-1)\cdot(-6)=-600+6=-594$.\par
\apantisi $-40$ και $-594$.\par
\elegxos $99\cdot6=594$ και ο αριθμός είναι αρνητικός.

\paradeigma{Παράσταση}
\tixero Η παράσταση $3\cdot(-4)-(-5)\cdot2+(-2)\cdot(-3)$.\par
\tithelo Να την υπολογίσω.\par
\possindeo Πρώτα οι πολλαπλασιασμοί: $3\cdot(-4)=-12$, $(-5)\cdot2=-10$, $(-2)(-3)=6$. Άρα $-12-(-10)+6=-12+10+6=4$.\par
\apantisi $4$.\par
\elegxos Προτεραιότητα: πολλαπλασιασμοί πριν από προσθέσεις και αφαιρέσεις.

\tmima{Ασκήσεις}

\zoni{Το κατάλαβα}{πρόσημα και τιμές}

\ex Συμπλήρωσε το μοτίβο: $3\cdot(-4)$, $2\cdot(-4)$, $1\cdot(-4)$, $0\cdot(-4)$, $(-1)\cdot(-4)$, $(-2)\cdot(-4)$, $(-3)\cdot(-4)$.

\ex Υπολόγισε: α) $(-6)\cdot(+4)$ \quad β) $(+7)\cdot(-3)$ \quad γ) $(-5)\cdot(-9)$ \quad δ) $(-12)\cdot0$ \quad ε) $(+8)\cdot(+5)$ \quad στ) $(-1)\cdot(-1)$.

\ex Υπολόγισε: α) $(+0{,}5)\cdot(-12)$ \quad β) $(-1{,}2)\cdot(-3)$ \quad γ) $(-2{,}5)\cdot(+0{,}4)$ \quad δ) $\left(-\frac{3}{4}\right)\cdot\left(-\frac{8}{9}\right)$ \quad ε) $\left(+\frac{2}{5}\right)\cdot(-15)$ \quad στ) $\left(-\frac{1}{2}\right)\cdot\left(-\frac{2}{3}\right)\cdot(-3)$.

\ex Βρες το πρόσημο του γινομένου χωρίς να το υπολογίσεις: α) $(-2)\cdot3\cdot(-5)\cdot(-1)$ \quad β) $(-1)(-1)(-1)(-1)$ \quad γ) $7\cdot(-3)\cdot(-2)\cdot4\cdot(-1)$ \quad δ) $(-5)\cdot0\cdot(-8)$.

\ex Υπολόγισε: α) $(-8)\cdot1$ \quad β) $(-8)\cdot(-1)$ \quad γ) $0\cdot(-3)$ \quad δ) $\alpha\cdot(-1)$ για $\alpha=-7$.

\ex Με την επιμεριστική: α) $(-4)\cdot(7+3)$ \quad β) $5\cdot(-3+8)$ \quad γ) $(-2)\cdot(5-9)$ \quad δ) $99\cdot(-6)$ \quad ε) $(-12)\cdot\left(\frac{1}{3}+\frac{1}{4}\right)$.

\ex Υπολόγισε: α) $3\cdot(-4)-(-5)\cdot2+(-2)\cdot(-3)$ \quad β) $(-2)\cdot(4-7)+3\cdot(-1-2)$ \quad γ) $10-2\cdot(-3)\cdot(-5)$.

\ex Σωστό ή λάθος; α) $(-3)(-4)=-12$. \quad β) Το γινόμενο ετεροσήμων είναι αρνητικό. \quad γ) $\alpha\cdot0=0$. \quad δ) $(-1)\cdot\alpha=-\alpha$. \quad ε) Το γινόμενο τριών αρνητικών είναι θετικό. \quad στ) $(-2)(-3)(-4)=-24$.

\zoni{Το χρησιμοποιώ}{θερμοκρασίες, χρέη, βουτιές}

\ex \textbf{Η θερμοκρασία.} Πέφτει $2$ βαθμούς κάθε ώρα. Πόσο θα έχει μεταβληθεί σε $3$ ώρες και πόσο έχει μεταβληθεί από πριν $3$ ώρες;

\ex \textbf{Το χρέος.} Κάθε εβδομάδα το χρέος της Μυρτώς αυξάνεται κατά $4$ €, δηλαδή η μεταβολή είναι $-4$ €. Ποια είναι η μεταβολή σε $5$ εβδομάδες;

\ex \textbf{Ο δύτης.} Ο δύτης κατεβαίνει $3$ μέτρα το λεπτό. Πού θα είναι σε $8$ λεπτά; Πού ήταν πριν από $5$ λεπτά;

\ex \textbf{Ο ανελκυστήρας.} Ο ανελκυστήρας κατεβαίνει $2$ ορόφους κάθε φορά, $4$ φορές. Πόση είναι η συνολική μεταβολή;

\ex \textbf{Το ταμείο.} Αν τα $9$ παιδιά της παρέας χρωστάνε από $4$ € το καθένα, τι αριθμός εκφράζει το συνολικό χρέος;

\ex \textbf{Η μετοχή.} Η τιμή μιας μετοχής πέφτει $0{,}25$ € κάθε ημέρα. Πόση είναι η συνολική μεταβολή σε $6$ ημέρες;

\zoni{Το σκέφτομαι}{γιατί το μείον επί μείον κάνει συν}

\ex Εξήγησε με το μοτίβο ή με τη θερμοκρασία γιατί $(-2)\cdot(-3)=+6$.

\ex Γιατί το τετράγωνο ενός ρητού δεν είναι ποτέ αρνητικό;

\ex Πόσα αρνητικά πρέπει να έχει ένα γινόμενο $10$ παραγόντων για να είναι θετικό; Δώσε δύο διαφορετικά πλήθη.

\ex Αν $\alpha\cdot\beta=0$, τι συμπεραίνεις για τους $\alpha,\,\beta$;

\ex Υπολόγισε το $(-1)\cdot(-1)\cdot\ldots\cdot(-1)$ με $100$ παράγοντες και με $101$ παράγοντες.

\zoni{Για το σπίτι}{φτιάξε το μοτίβο}

\ex Φτιάξε το μοτίβο του πολλαπλασιασμού με το $-5$ από $3$ ως $-3$ και πες τι παρατηρείς.
